\section{Discussion}

The proposed framework investigates whether prompt engineering can improve
LLM-assisted algorithmic problem solving and potentially change the role of
traditional algorithm practice.

The central motivation of this study is not to suggest that traditional data
structures and algorithms knowledge will become unnecessary. Instead, the
research examines whether the skills required for solving algorithmic problems
may shift as Large Language Models become increasingly capable of generating
and optimizing code.

\subsection{Traditional Algorithmic Skills}

Traditional LeetCode-style preparation develops several important skills,
including algorithm recognition, data structure selection, complexity analysis,
and independent implementation.

These skills remain important even when AI tools are available. Developers must
still understand whether an AI-generated solution is correct, whether its
complexity is appropriate, and whether it satisfies the requirements of a
given problem.

Therefore, AI-assisted programming should not necessarily be interpreted as a
replacement for algorithmic knowledge.

\subsection{The Emerging Role of Prompt Engineering}

The increasing capability of LLMs introduces an additional skill: the ability
to communicate programming requirements effectively to an AI system.

A developer may use structured prompts to request algorithm comparison,
complexity analysis, edge-case verification, and code optimization.

This suggests a possible shift from a purely implementation-oriented workflow
toward a collaborative workflow involving problem formulation, AI interaction,
and solution verification.

\subsection{Implications for Computer Science Education}

If future experiments demonstrate that structured prompting substantially
improves algorithmic problem-solving performance, computer science education
may need to consider how traditional algorithm practice and AI-assisted
programming should be combined.

Rather than eliminating data structures and algorithms education, AI tools may
change how these concepts are practiced.

Students may increasingly need both foundational algorithmic knowledge and the
ability to evaluate and guide AI-generated solutions.

\subsection{Implications for Software Engineering Interviews}

Technical interviews traditionally use algorithmic problems to evaluate a
candidate's ability to reason about data structures, algorithms, and
complexity.

The increasing availability of AI programming assistants raises questions
about whether these evaluation methods will remain representative of real-world
software engineering workflows.

If developers routinely use AI assistants in professional environments,
future assessments may place greater emphasis on problem formulation,
architecture, debugging, code review, and AI-assisted development.

However, these implications require empirical validation and should not be
interpreted as evidence that traditional algorithm interviews are obsolete.

\subsection{Interpretation of the Research Question}

The research question should therefore be interpreted as an investigation of
skill transformation rather than complete skill replacement.

The potential transition can be represented as:

\[
\text{Algorithm Practice}
\rightarrow
\text{AI-Assisted Algorithm Problem Solving}
\]

rather than:

\[
\text{Algorithm Practice}
\rightarrow
\text{No Algorithm Knowledge Required}
\]

The proposed benchmark provides a framework for empirically investigating this
transition.
